\documentclass[10pt,a4paper]{amsart}
\usepackage[margin=19mm]{geometry}
\usepackage{amsmath,amssymb,mathtools}
\usepackage[scr=rsfs,cal=euler]{mathalfa}
\usepackage{xfrac}
\usepackage[unicode,hidelinks,bookmarks=false]{hyperref}
\newcommand{\ie}{i.e.\ }
\newcommand{\cf}{cf.\ }
\newcommand{\resp}{respectively\ }
\newcommand{\Subsection}{\subsection}
\providecommand{\nobreakdash}{\nobreak\hbox{-}}
\setlength{\emergencystretch}{2em}
\multlinegap0pt
\usepackage[shortlabels]{enumitem}
\usepackage{mathtools}
\usepackage{bm}
\usepackage{amssymb}
\usepackage{dsfont}
\usepackage{float}
\newtheorem{proposition}{Proposition}
\newtheorem{definition}{Definition}
\newtheorem{lemma}{Lemma}
\newcommand{\Law}{\mathrm{Law}}
\newcommand{\TAP}{\mathrm{TAP}}
\newcommand{\bfm}{\mathbf{m}}
\newcommand{\dE}{\mathbb{E}}%
\newcommand{\dP}{\mathbb{P}}%
\newcommand{\dR}{\mathbb{R}}%
\newcommand{\dd}{\mathrm{d}}
\newcommand{\eq}{\mathrm{eq}}
\newcommand{\mc}{\mathcal}
\newcommand{\ve}{\varepsilon}
\title{Large deviations for the maximum of\\ the generalized TAP free energy}
\author{Jeanne Boursier}
\date{Source: arXiv:2607.29330v1 (CC BY 4.0)}
\begin{document}
\maketitle

\noindent\emph{Source and licence.} Large deviations for the maximum of\\ the generalized TAP free energy. Version arXiv:2607.29330v1 (CC BY 4.0), distributed by arXiv under the Creative Commons Attribution 4.0 International licence, http://creativecommons.org/licenses/by/4.0/, as recorded in the arXiv licence metadata for this exact version and shown on the abstract page https://arxiv.org/abs/2607.29330v1. Full licence notice: \texttt{licenses/arxiv-2607.29330-CC-BY-4.0.txt}.\par\smallskip

\noindent\textit{Editorial scope.} Counted unit: the article's proposition prop:onshell-lower-bound (source0.tex lines 2579--2598). The article's complete original proof of it is delivered with it (source0.tex lines 3638--4042, one \texttt{proof} environment of 405 lines, which the article places away from the statement). No mathematics is written, repaired or re-derived: the statement and the proof are the article's own lines, in the article's own notation, and no retained line is substituted or re-flowed.  Retained uncounted as the source-local context the proof needs: lines 2606--2635; lines 2644--2673; lines 2712--2739; lines 3102--3153; lines 3525--3534; lines 3568--3582; lines 4594--4662; lines 4632--4643; lines 4671--4737.  Nothing was omitted from any retained span, so the package carries no omission record.  No retained line contains a citation, so the wrapper carries no bibliography and no bibliography entry.  No retained line contains a figure environment, an image-inclusion command or drawing code, so no image is delivered and the package declares no asset dependency.  Editorial formatting only, in addition: an AMS \texttt{amsart} preamble that restates the article's own declarations and macros used by the retained lines, namely the environments \texttt{proposition}, \texttt{definition}, \texttt{lemma} on separate counters, together with the standard packages those lines need.  The retained environments print in this document as Proposition 1, Definition 1, Lemma 1 in order of appearance, whereas the article prints the same items with the numbering of its own full document.  The complete native source of the article as distributed in the arXiv e-print for this exact version is bundled unchanged as \texttt{sources/206474/source0.tex}.
	\begin{definition}[SUSY states]\label{def:susy-state}
		Fix a regular \(f\in(f_{\eq},f_1)\), write \(\theta=\theta_f\) and
		\(\zeta=\zeta_f\), and let \(q\) be the first positive contact point of
		\(\zeta\). Let \(X\) be the Auffinger--Chen process associated with
		\(\zeta\).
		
		For \(\ve>0\), an \(\ve\)-SUSY state at level \(f\) is a
		critical point \(\bfm\in(-1,1)^N\) of \(F_{\TAP,\zeta}\) satisfying
		\[
		|q_\bfm-q|\le\ve,
		\qquad
		\left|\frac1N F_{\TAP,\zeta}(\bfm)-f\right|\le\ve,
		\qquad
		\left|\frac1N F_{\TAP}(\bfm)-f\right|\le\ve.
		\]
		Define \(b_i\) by
		\[
		m_i=\partial_x\Phi_\zeta(q,b_i),
		\qquad 1\le i\le N,
		\]
		and require
		\[
		W_2\left(
		\frac1N\sum_{i=1}^N\delta_{b_i},
		\Law(X_q)
		\right)\le\ve.
		\]
		Let \(\mathsf N_\ve\) denote the number of
		\(\ve\)-SUSY states at level \(f\).
	\end{definition}

	\begin{definition}[Isolated SUSY states]\label{def:isolated-susy-state}
		For \(\ve,\kappa>0\) and \(0<\delta<q/4\), an
		\(\ve\)-SUSY state \(\bfm\) is isolated at parameters
		\((\delta,\kappa)\) if:
		\begin{enumerate}[label=\textup{(\roman*)}]
			\item no distinct \(\ve\)-SUSY state \(\bfm'\)
			satisfies
			\[
			R(\bfm,\bfm')\in[\delta,q-\delta].
			\]
			\item
			\[
			\sup_{\substack{
					x\in(-1,1)^N,\\
					\|x-\bfm\|_2\le2\sqrt{\delta N}
			}}
			\lambda_{\max}\!\left(
			\nabla^2F_{\TAP,\zeta}(x)
			\right)
			\le-\kappa.
			\]
		\end{enumerate}
		Let
		\[
		\widehat{\mathsf N}_{\ve,\delta,\kappa}
		:=
		\#\{\text{\(\ve\)-SUSY states isolated at parameters
			\((\delta,\kappa)\)}\}.
		\]
	\end{definition}

	\begin{proposition}[Annealed complexity of SUSY states]\label{prop:susy-annealed}
		Let \(f\in(f_{\eq},f_1)\) be regular, write
		\(\theta=\theta_f\) and \(\zeta=\zeta_f\), and let \(q\) be the first
		positive contact point of \(\zeta\). Let \(X\) be the
		Auffinger--Chen process associated with \(\zeta\). Assume that
		\[
		1-\xi''(q)\dE\left[
		\left(
		\partial_{xx}\Phi_\zeta(q,X_q)
		\right)^2
		\right]>0.
		\]
		Then the counts \(\mathsf N_\ve\) from
		Definition~\ref{def:susy-state} satisfy
		\[
		\liminf_{\ve\downarrow0}
		\liminf_{N\to\infty}
		\frac1N\log\dE\mathsf N_\ve
		=
		\limsup_{\ve\downarrow0}
		\limsup_{N\to\infty}
		\frac1N\log\dE\mathsf N_\ve
		=
		\theta\bigl(P_\zeta-f\bigr)
		=
		\Sigma(f).
		\]
	\end{proposition}

	\begin{proposition}\label{prop:sibling-guerra-bound}
		Fix a regular \(f\in(f_{\eq},f_1)\), write
		\(\theta=\theta_f\) and \(\zeta=\zeta_f\), and let \(q\) be the first
		positive contact point of \(\zeta\). Let \(X\) be the
		Auffinger--Chen process associated with \(\zeta\).
		For \(\ve>0\), fix \(\bfm\in(-1,1)^N\), and define \(b_i\) by
		\[
		m_i=\partial_x\Phi_\zeta(q,b_i).
		\]
		Assume that
		\[
		|q_\bfm-q|\le\ve,
		\qquad
		W_2\left(
		\frac1N\sum_{i=1}^N\delta_{b_i},
		\Law(X_q)
		\right)\le\ve.
		\]
		Let \(f'\in[f-\ve,f+\ve]\) satisfy
		\[
		\left|
		f'
		+\frac1N\left(
		F_{\TAP}(\bfm)-F_{\TAP,\zeta}(\bfm)
		\right)
		-f
		\right|
		\le\ve.
		\]
		Let \(\dE_{\bfm,f'}\) denote conditional expectation given
		\[
		F_{\TAP,\zeta}(\bfm)=Nf',
		\qquad
		\nabla F_{\TAP,\zeta}(\bfm)=0.
		\]
		For every \(r\in(0,q)\),
		every \(\delta>0\), and every fixed finite-step probability measure
		\(\rho\) on \([r,q]\),
		\begin{equation}\label{eq:sibling-guerra-bound}
			\frac1N\dE_{\bfm,f'}
			\max_{\bfm'\in\mc C_N(\bfm,r,\delta)}
			\left[
			F_{\TAP,\zeta}(\bfm')-F_{\TAP,\zeta}(\bfm)
			\right]
			\le
			\mc J_{N,\bfm,q,r}(\rho)
			+o_\ve(1)+o_\delta(1)+o_N(1).
		\end{equation}
		The errors are uniform over the displayed \(\ve\)-window data,
		with \(N\to\infty\) taken before
		\(\ve,\delta\downarrow0\).
	\end{proposition}

	\begin{equation}\label{eq:sibling-nonpositive-bound}
		\frac1N\dE_{\bfm,f'}
		\max_{\bfm'\in\mc C_N(\bfm,r,\delta)}
		\left[
		F_{\TAP,\zeta}(\bfm')-F_{\TAP,\zeta}(\bfm)
		\right]
		\le
		-\theta H_{\zeta,q}(r)
		+o_\ve(1)+o_\delta(1)+o_N(1),
	\end{equation}

	\begin{lemma}
		\label{lem:one-sided-spherical-code}
		For every \(a\in(0,1)\),
		\begin{equation}\label{eq:one-sided-spherical-code-explicit}
			\limsup_{N\to\infty}
			\frac1N\log A_N(a)
			\le-\frac12\log(1-a).
		\end{equation}
		In particular,
		\begin{equation}\label{eq:one-sided-spherical-code}
			\lim_{a\downarrow0}
			\limsup_{N\to\infty}
			\frac1N\log A_N(a)=0.
		\end{equation}
	\end{lemma}

	\begin{proposition}\label{prop:one-point-conditional-hessian}
		Let \(f\in(f_{\eq},f_1)\) be regular, let \(\zeta=\zeta_f\), and let
		\(q\) be the first positive contact point of \(\zeta\). Let \(X\) be the
		corresponding Auffinger--Chen process. Suppose that
		\[
		1-\xi''(q)\dE\left[
		\left(
		\partial_{xx}\Phi_\zeta(q,X_q)
		\right)^2
		\right]>0.
		\]
		Let \(\bfm\in (-1,1)^N\) satisfy \(q_\bfm=q\), and define \(b_i\) by
		\(m_i=\partial_x\Phi_\zeta(q,b_i)\). Assume that
		\[
		\frac1N\sum_{i=1}^N\delta_{b_i}
		\xrightarrow{W_2}\Law(X_q).
		\]
		Let \(f_N\to f\) and condition on
		\[
		\nabla F_{\TAP,\zeta}(\bfm)=0,
		\qquad
		F_{\TAP,\zeta}(\bfm)=Nf_N.
		\]
		Write \(\dP_{\bfm,f_N}\) and \(\dE_{\bfm,f_N}\) for the corresponding
		conditional probability and expectation. For
		\(\sigma\in\{-,+\}\), the following conclusions hold.
		\begin{enumerate}[label=\textup{(\roman*)}]
			\item There is \(c>0\) such that
			\begin{equation}
				\dP_{\bfm,f_N}\left(
				\nabla_\sigma^2F_{\TAP,\zeta}(\bfm)\preceq-cI_N
				\right)\longrightarrow1.
			\end{equation}
			
			\item
			\begingroup
			There are \(\delta_0,\kappa>0\) such that, for every
			\(0<\delta<\delta_0\),
			\begin{equation}\label{eq:one-point-local-spectral-gap}
				\dP_{\bfm,f_N}\left(
				\sup_{\substack{
						x\in(-1,1)^N,\\
						\|x-\bfm\|_2\le2\sqrt{\delta N}
				}}
				\lambda_{\max}\!\left(
				\nabla^2F_{\TAP,\zeta}(x)
				\right)
				\le-\kappa
				\right)\longrightarrow1.
			\end{equation}
			If \(q_x\) is an atom of \(\zeta\), the two one-sided Hessians may
			differ. In \eqref{eq:one-point-local-spectral-gap},
			\(\nabla^2F_{\TAP,\zeta}(x)\) denotes
			\(\nabla_-^2F_{\TAP,\zeta}(x)\), obtained by approaching \(q_x\)
			through smaller self-overlaps.
			\endgroup
		\end{enumerate}
		The conclusion in \textup{(ii)} is uniform under
		\[
		|q_\bfm-q|\le\ve,
		\qquad
		|f_N-f|\le\ve,
		\qquad
		W_2\left(N^{-1}\sum_{i=1}^N\delta_{b_i},\Law(X_q)\right)\le\ve.
		\]
		The probability of the complementary event in
		\eqref{eq:one-point-local-spectral-gap} is
		\(o_\ve(1)+o_N(1)\), with \(N\to\infty\) before \(\ve\downarrow0\).
	\end{proposition}

			\begin{equation}
				\dP_{\bfm,f_N}\left(
				\sup_{\substack{
						x\in(-1,1)^N,\\
						\|x-\bfm\|_2\le2\sqrt{\delta N}
				}}
				\lambda_{\max}\!\left(
				\nabla^2F_{\TAP,\zeta}(x)
				\right)
				\le-\kappa
				\right)\longrightarrow1.
			\end{equation}

	\begin{lemma}[Conditional determinant estimates]
		\label{lem:one-point-conditional-determinant}
		Under the hypotheses and notation of
		Proposition~\ref{prop:one-point-conditional-hessian}, the following
		hold for \(\sigma\in\{-,+\}\):
		\begin{enumerate}[label=\textup{(\roman*)}]
			\item
			\begin{equation}\label{eq:one-point-conditional-determinant}
				\begin{aligned}
					\frac1N
					\log\dE_{\bfm,f_N}
					\left|\det\nabla_\sigma^2F_{\TAP,\zeta}(\bfm)\right|
					={}&
					-\frac1N\sum_{i=1}^N
					\log\partial_{xx}\Phi_\zeta(q,b_i)
					+\frac{\xi''(q)}2
					\left(
					\frac1N\sum_{i=1}^N
					\partial_{xx}\Phi_\zeta(q,b_i)
					\right)^2
					+o(1)\\
					\longrightarrow{}&
					-\dE\log\partial_{xx}\Phi_\zeta(q,X_q)
					+\frac{\xi''(q)}2\widehat\zeta(q)^2 .
				\end{aligned}
			\end{equation}
			In the \(\ve\)-window above, the error in
			\eqref{eq:one-point-conditional-determinant} is
			\(o_\ve(1)+o_N(1)\).
			
			\item
			Suppose that, for some \(\kappa_0>0\) and \(L<\infty\),
			\[
			\frac1N\sum_{i=1}^N
			\exp\{(8+\kappa_0)|b_i|\}\le L.
			\]
			Set
			\[
			D_{\bfm,\sigma}
			:=
			\det\nabla_\sigma^2F_{\TAP,\zeta}(\bfm).
			\]
			There are \(c_L>0\) and a deterministic \(r_N=o(N)\) such that
			\begin{equation}\label{eq:one-point-determinant-lower-tail}
				\dP_{\bfm,f_N}\left(
				|D_{\bfm,\sigma}|
				<e^{-r_N}\dE_{\bfm,f_N}|D_{\bfm,\sigma}|
				\right)
				\le e^{-c_LN}.
			\end{equation}
			Moreover, the estimate is uniform over the same window, subject to the
			exponential-moment bound above. There is a function \(\rho_{\det}\),
			with \(\rho_{\det}(\ve)\downarrow0\), and there are nonnegative
			deterministic numbers \(r_{N,\ve}=o_N(N)\) such that
			\begin{equation}\label{eq:susy-one-point-determinant-lower-tail}
				\dP_{\bfm,f_N}\left(
				|D_{\bfm,\sigma}|
				<
				\exp\{-N\rho_{\det}(\ve)-r_{N,\ve}\}
				\dE_{\bfm,f_N}|D_{\bfm,\sigma}|
				\right)
				\le e^{-c_LN}.
			\end{equation}
			Here \(r_{N,\ve}/N\to0\) for fixed \(\ve\), and \(N\to\infty\)
			before \(\ve\downarrow0\).
		\end{enumerate}
	\end{lemma}

	\begin{proposition}\label{prop:onshell-lower-bound}
		Let \(f\in(f_{\eq},f_1)\) be regular. Let \(q\) be the first
		positive contact point of \(\zeta_f\), and let \(X\) be the Auffinger--Chen process
		associated with \(\zeta_f\).
		Assume that
		\[
		1-\xi''(q)\dE\left[
		\left(
		\partial_{xx}\Phi_{\zeta_f}(q,X_q)
		\right)^2
		\right]>0.
		\]
		Then
		\[
		\lim_{\ve\downarrow0}
		\liminf_{N\to\infty}
		\frac1N\log\dP\bigl(\mc E_{N,\ve}(f)\bigr)
		\ge\Sigma(f).
		\]
	\end{proposition}

	\begin{proof}[Proof of Proposition~\ref{prop:onshell-lower-bound}]
		Write \(\theta=\theta_f\) and \(\zeta=\zeta_f\).
		Let \(\mathsf N_\ve\) be the SUSY count from
		Definition~\ref{def:susy-state}.
		Proposition~\ref{prop:one-point-conditional-hessian} gives
		\(\delta_0,\kappa>0\). Fix
		\[
		0<\delta<\min\left\{\frac q4,\delta_0\right\}.
		\]
		Fix \(\kappa_0>0\) and choose
		\[
		L>
		\dE\exp\{(8+\kappa_0)|X_q|\}.
		\]
		Let \(\ve_{\mathrm H}(\delta)>0\) and
		\(\ve_{\det}(L)>0\) be thresholds for the
		local spectral-gap and determinant lower-tail estimates.
		
		For a deterministic \(\bfm\) and \(f'\in\dR\), let
		\(\dP_{\bfm,f'}\) and \(\dE_{\bfm,f'}\) denote probability and
		expectation conditioned on
		\[
		F_{\TAP,\zeta}(\bfm)=Nf',\quad
		\nabla F_{\TAP,\zeta}(\bfm)=0.
		\]
		
		\medskip\noindent\emph{Step 1. Stable truncation.}
		The regression and comparison errors are locally uniform in
		\(r\in[\delta,q-\delta]\). Indeed, the covariance coefficients,
		deterministic terms, and reduced Parisi equation in
		Proposition~\ref{prop:sibling-guerra-bound} depend continuously on \(r\)
		and remain uniformly bounded on this interval.
		For each \(r\in[\delta,q-\delta]\),
		\eqref{eq:sibling-nonpositive-bound} and
		\(\theta H_{\zeta,q}(r)\ge8g_\delta\) therefore give
		\(h_r,\ve_r>0\) such that, for all sufficiently large \(N\),
		\[
		\frac1N\dE_{\bfm,f'}
		\max_{\bfm'\in\mc C_N(\bfm,r,h_r)}
		\left[
		F_{\TAP}(\bfm')-F_{\TAP}(\bfm)
		\right]
		\le-6g_\delta
		\]
		uniformly over the one-point windows of width at most
		\(\ve_r\).
		
		Choose \(r_1,\ldots,r_J\) such that the intervals
		\[
		\left(r_j-\frac{h_{r_j}}2,r_j+\frac{h_{r_j}}2\right),
		\qquad 1\le j\le J,
		\]
		cover \([\delta,q-\delta]\). Choose
		\[
		0<\ve_{\mathrm I}
		\le
		\min_{1\le j\le J}
		\left\{\ve_{r_j},h_{r_j}\right\}.
		\]
		If \(0<\ve<\ve_{\mathrm I}\), every point in the maximum
		below belongs to at least one slice
		\(\mc C_N(\bfm,r_j,h_{r_j})\), and the expected maximum on each slice is at
		most \(-6Ng_\delta\). The conditional variance of every Hamiltonian
		increment on the cube is bounded by \(CN\). Gaussian concentration and
		a union bound over the finite cover give \(c_\delta^{\mathrm I}>0\) such
		that
		\begin{equation}\label{eq:sibling-uniform-intermediate}
			\dP_{\bfm,f'}\left(
			\max_{\substack{
					\bfm'\in[-1,1]^N,\\
					|q_{\bfm'}-q|\le\ve,\\
					R(\bfm,\bfm')\in[\delta,q-\delta]
			}}
			\left[F_{\TAP}(\bfm')-F_{\TAP}(\bfm)\right]
			>-2Ng_\delta
			\right)
			\le e^{-c_\delta^{\mathrm I}N}
		\end{equation}
		for \(0<\ve<\ve_{\mathrm I}\), uniformly over the
		one-point Kac--Rice windows.
		
		Choose \(\ve>0\) such that
		\begin{equation}\label{eq:susy-geometric-parameter-choice}
			\ve
			<
			\min\left\{
			\ve_{\mathrm H}(\delta),\ve_{\det}(L),
			\ve_{\mathrm I},
			\delta,g_\delta
			\right\}.
		\end{equation}
		
		For a point \(\bfm\) counted by \(\mathsf N_\ve\), let
		\(\mathcal G_{\bfm,\delta}\) be the event that \(\bfm\) is isolated in
		the sense of Definition~\ref{def:isolated-susy-state}.
		
		If both \(\bfm\) and \(\bfm'\) are counted by
		\(\mathsf N_\ve\), their TAP values differ by at most
		\(2N\ve\). Hence, by
		\eqref{eq:susy-geometric-parameter-choice}, failure of condition
		\textup{(i)} in Definition~\ref{def:isolated-susy-state} is contained
		in the event in \eqref{eq:sibling-uniform-intermediate}. Therefore
		\begin{equation}\label{eq:geometric-intermediate-failure}
			\dP_{\bfm,f'}\left(
			\bfm\text{ fails condition \textup{(i)} in
				Definition~\ref{def:isolated-susy-state}}
			\right)
			\le e^{-c_\delta^{\mathrm I}N}.
		\end{equation}
		By \eqref{eq:one-point-local-spectral-gap}, for all sufficiently
		large \(N\), we use the local spectral-gap estimate only in the form
		\begin{equation}\label{eq:geometric-local-gap-probability}
			\dP_{\bfm,f'}\left(
			\bfm\text{ satisfies condition \textup{(ii)} in
				Definition~\ref{def:isolated-susy-state}}
			\right)
			\ge\frac12.
		\end{equation}
		\eqref{eq:geometric-intermediate-failure} and
		\eqref{eq:geometric-local-gap-probability} give
		\begin{equation}\label{eq:geometric-good-point-probability}
			\dP_{\bfm,f'}\left(
			\mathcal G_{\bfm,\delta}
			\right)
			\ge\frac12-e^{-c_\delta^{\mathrm I}N}.
		\end{equation}
		
		For \(f'\in[f-\ve,f+\ve]\), let
		\[
		\mathcal W_{\ve,f'}
		:=
		\left\{
		\begin{array}{l}
			\bfm\in(-1,1)^N:\ |q_\bfm-q|\le\ve,\\[1mm]
			W_2\left(N^{-1}\sum_{i=1}^N\delta_{b_i},\Law(X_q)\right)
			\le\ve,\\[1mm]
			\left|f'+N^{-1}\left(F_{\TAP}(\bfm)
			-F_{\TAP,\zeta}(\bfm)\right)-f\right|\le\ve
		\end{array}
		\right\}.
		\]
		Set
		\[
		\mathcal W_{\ve,f'}^{(L)}
		:=
		\left\{
		\bfm\in\mathcal W_{\ve,f'}:
		\frac1N\sum_{i=1}^N
		\exp\{(8+\kappa_0)|b_i|\}\le L
		\right\}.
		\]
		By the Kac--Rice formula,
		\[
		\begin{aligned}
			\dE\widehat{\mathsf N}_{\ve,\delta,\kappa}
			\ge{}&
			N\int_{f-\ve}^{f+\ve}
			\int_{\mathcal W_{\ve,f'}^{(L)}}\\
			&\quad\dE_{\bfm,f'}\!\left[
			\left|\det\nabla^2F_{\TAP,\zeta}(\bfm)\right|
			\boldsymbol 1_{\mathcal G_{\bfm,\delta}}
			\right]\\
			&\quad\times
			p_{\left(
				\nabla F_{\TAP,\zeta}(\bfm),
				F_{\TAP,\zeta}(\bfm)
				\right)}(0,Nf')
			\,\dd\bfm\,\dd f'.
		\end{aligned}
		\]
		
		Set
		\[
		D_\bfm
		:=
		\det\nabla^2F_{\TAP,\zeta}(\bfm).
		\]
		Let \(r_{N,\ve}=o_N(N)\) be the uniform remainder from
		Lemma~\ref{lem:one-point-conditional-determinant}\textup{(ii)}, and let
		\[
		\mathcal D_\bfm
		:=
		\left\{
		|D_\bfm|
		\ge
		\exp\{-N\rho_{\det}(\ve)-r_{N,\ve}\}
		\dE_{\bfm,f'}|D_\bfm|
		\right\}.
		\]
		Lemma~\ref{lem:one-point-conditional-determinant}\textup{(ii)} gives
		\[
		\dP_{\bfm,f'}(\mathcal D_\bfm^{\,c})
		\le e^{-c_LN}
		\]
		uniformly on \(\mathcal W_{\ve,f'}^{(L)}\). Together with
		\eqref{eq:geometric-good-point-probability}, this gives
		\[
		\dP_{\bfm,f'}\left(
		\mathcal G_{\bfm,\delta}\cap\mathcal D_\bfm
		\right)
		\ge
		\frac12-e^{-c_\delta^{\mathrm I}N}-e^{-c_LN}
		\ge\frac14
		\]
		for all sufficiently large \(N\). Therefore
		\begin{equation}\label{eq:geometric-local-stability-input}
			\dE_{\bfm,f'}\left[
			|D_\bfm|
			\boldsymbol 1_{\mathcal G_{\bfm,\delta}}
			\right]
			\ge
			\frac14
			\exp\{-N\rho_{\det}(\ve)-r_{N,\ve}\}
			\dE_{\bfm,f'}|D_\bfm|
		\end{equation}
		uniformly on \(\mathcal W_{\ve,f'}^{(L)}\).
		
		The Auffinger--Chen drift is bounded, so \(X_q\) has all exponential
		moments. With \(L\) chosen above,
		the law of large numbers shows that the cutoff has probability tending to
		one under the product measure used in the lower-bound proof of
		Proposition~\ref{prop:susy-annealed}. Hence the restricted
		Kac--Rice integral has exponent \(\Sigma(f)\). Together with
		\eqref{eq:geometric-local-stability-input} and
		\(\rho_{\det}(\ve)\to0\), this gives the lower bound. The upper
		bound follows from
		\(\widehat{\mathsf N}_{\ve,\delta,\kappa}\le\mathsf N_\ve\) and
		Proposition~\ref{prop:susy-annealed}. Therefore
		\begin{equation}\label{eq:geometric-truncated-first-moment}
			\lim_{\ve\downarrow0}
			\liminf_{N\to\infty}
			\frac1N
			\log
			\dE\widehat{\mathsf N}_{\ve,\delta,\kappa}
			=
			\Sigma(f).
		\end{equation}
		
		\medskip\noindent\emph{Step 2. Exclusion at high overlap.}
		A point \(\bfm\) counted by
		\(\widehat{\mathsf N}_{\ve,\delta,\kappa}\) has no distinct competitor
		\(\bfm'\) counted by \(\mathsf N_\ve\) with
		\[
		R(\bfm,\bfm')\ge q-\delta.
		\]
		To prove this, suppose that such a competitor exists. Since
		\(\ve\le\delta\),
		\[
		\begin{aligned}
			\frac1N\|\bfm'-\bfm\|_2^2
			&=
			q_\bfm+q_{\bfm'}-2R(\bfm,\bfm')\\
			&\le
			2(q+\ve)-2(q-\delta)
			\le4\delta.
		\end{aligned}
		\]
		With \(v=\bfm'-\bfm\), the segment
		\[
		\{\bfm+tv:0\le t\le1\}
		\]
		is contained in the ball appearing in condition \textup{(ii)} of
		Definition~\ref{def:isolated-susy-state}.
		
		Since both endpoints are critical,
		\[
		\nabla F_{\TAP,\zeta}(\bfm)
		=
		\nabla F_{\TAP,\zeta}(\bfm')
		=0.
		\]
		Taylor's formula for the gradient gives
		\[
		\begin{aligned}
			0
			&=
			\left\langle
			v,
			\nabla F_{\TAP,\zeta}(\bfm')
			-
			\nabla F_{\TAP,\zeta}(\bfm)
			\right\rangle\\
			&=
			\int_0^1
			v^{\mathsf T}
			\nabla^2F_{\TAP,\zeta}(\bfm+tv)
			v\,\dd t\\
			&\le
			-\kappa\|v\|_2^2.
		\end{aligned}
		\]
		It follows that \(v=0\), contradicting \(\bfm'\ne\bfm\).
		
		\medskip\noindent\emph{Step 3. Spherical-code bound.}
		Let \(\bfm\) and \(\bfm'\) be two distinct points counted by
		\(\widehat{\mathsf N}_{\ve,\delta,\kappa}\). By condition
		\textup{(i)} of Definition~\ref{def:isolated-susy-state}, their overlap
		cannot belong to
		\([\delta,q-\delta]\), and Step~2 excludes overlaps at least
		\(q-\delta\). Therefore
		\[
		R(\bfm,\bfm')<\delta.
		\]
		
		Normalize the retained points by setting
		\[
		u_\bfm:=\frac{\bfm}{\sqrt{q_\bfm}}.
		\]
		Then
		\[
		\|u_\bfm\|_2=\sqrt N,
		\]
		and, for two distinct retained points,
		\[
		\frac1N
		\langle u_\bfm,u_{\bfm'}\rangle
		=
		\frac{R(\bfm,\bfm')}
		{\sqrt{q_\bfm q_{\bfm'}}}
		\le
		\frac{\delta}{q-\ve}.
		\]
		Thus the normalized retained points form a code to which
		Lemma~\ref{lem:one-sided-spherical-code} applies, and
		\begin{equation}\label{eq:geometric-code-count}
			\widehat{\mathsf N}_{\ve,\delta,\kappa}
			\le
			A_N\left(
			\frac{\delta}{q-\ve}
			\right).
		\end{equation}
		
		\medskip\noindent\emph{Step 4. Existence probability.}
		By \eqref{eq:geometric-code-count},
		\[
		\widehat{\mathsf N}_{\ve,\delta,\kappa}
		\le
		A_N\left(
		\frac{\delta}{q-\ve}
		\right)
		\boldsymbol 1_{
			\{\widehat{\mathsf N}_{\ve,\delta,\kappa}\ge1\}
		}.
		\]
		Taking expectations gives
		\[
		\dP\left(
		\widehat{\mathsf N}_{\ve,\delta,\kappa}\ge1
		\right)
		\ge
		\frac{
			\dE\widehat{\mathsf N}_{\ve,\delta,\kappa}
		}{
			A_N\left(\delta/(q-\ve)\right)
		}.
		\]
		Since
		\(\mathsf N_\ve
		\ge\widehat{\mathsf N}_{\ve,\delta,\kappa}\),
		\[
		\dP\bigl(\mc E_{N,\ve}(f)\bigr)
		\ge
		\dP\left(
		\widehat{\mathsf N}_{\ve,\delta,\kappa}\ge1
		\right).
		\]
		Lemma~\ref{lem:one-sided-spherical-code} now gives
		\[
		\begin{aligned}
			\lim_{\ve\downarrow0}
			\liminf_{N\to\infty}
			\frac1N
			\log
			\dP\bigl(\mc E_{N,\ve}(f)\bigr)
			&\ge
			\lim_{\ve\downarrow0}
			\left[
			\liminf_{N\to\infty}
			\frac1N
			\log\dE\widehat{\mathsf N}_{\ve,\delta,\kappa}
			+
			\frac12\log\left(
			1-\frac{\delta}{q-\ve}
			\right)
			\right]\\
			&=
			\Sigma(f)
			+
			\frac12\log\left(1-\frac{\delta}{q}\right),
		\end{aligned}
		\]
		where the equality follows from
		\eqref{eq:geometric-truncated-first-moment}. Letting \(\delta\downarrow0\)
		gives
		\[
		\lim_{\ve\downarrow0}
		\liminf_{N\to\infty}
		\frac1N
		\log
		\dP\bigl(\mc E_{N,\ve}(f)\bigr)
		\ge
		\Sigma(f),
		\]
		which proves Proposition~\ref{prop:onshell-lower-bound}.
	\end{proof}

\end{document}
